\item ¿Por qué es imposible la siguiente situación? Un astronauta en la Luna está studying movimiento ondulatorio mediante el aparato analizado en el ejemplo 16.3 y mostrado en la figura 16.12. Él mide el intervalo de tiempo empleado por los pulsos para viajar a lo largo del alambre horizontal. Suponga que el alambre horizontal tiene una masa de 4.00 g, una longitud de 1.60 m y que un objeto de 3.00 kg está suspendido de su extensión alrededor de la polea. El astronauta encuentra que un pulso requiere 26.1 ms para recorrer la longitud del alambre.